\documentclass[11pt,a4paper]{article}

% --------------------------------------------------
% Packages
% --------------------------------------------------
\usepackage[utf8]{inputenc}
\usepackage[T1]{fontenc}
\usepackage[french]{babel}
\usepackage[a4paper,margin=2cm]{geometry}

\usepackage{amsmath,amssymb,amsthm,mathtools}
\usepackage{lmodern}
\usepackage{microtype}
\usepackage{enumitem}
\usepackage{tikz}
\usepackage{pgfplots}
\usepackage{tcolorbox}
\usepackage{xcolor}
\usepackage{fancyhdr}
\usepackage{hyperref}
\usepackage{titlesec}
\usepackage{array}
\usepackage{booktabs}
\usepackage{multicol}
\usepackage{wrapfig}
\usepackage{graphicx}

\usetikzlibrary{arrows.meta,calc,decorations.pathreplacing,angles,quotes,positioning}
\pgfplotsset{compat=1.18}

% --------------------------------------------------
% Couleurs
% --------------------------------------------------
\definecolor{bleufonce}{RGB}{24,62,120}
\definecolor{bleuclair}{RGB}{230,239,252}
\definecolor{grisclair}{RGB}{245,245,245}
\definecolor{vertfonce}{RGB}{40,110,80}
\definecolor{rougefonce}{RGB}{150,50,50}

% --------------------------------------------------
% Hyperref
% --------------------------------------------------
\hypersetup{
    colorlinks=true,
    linkcolor=bleufonce,
    urlcolor=bleufonce,
    citecolor=bleufonce
}

% --------------------------------------------------
% Mise en page
% --------------------------------------------------
\setlength{\parindent}{0pt}
\setlength{\parskip}{0.45em}

\pagestyle{fancy}
\fancyhf{}
\fancyhead[L]{\textsc{DM de Physique-Chimie}}
\fancyhead[R]{\textsc{Première générale}}
\fancyfoot[C]{\thepage}

% --------------------------------------------------
% Titres
% --------------------------------------------------
\titleformat{\section}
{\Large\bfseries\color{bleufonce}}
{}{0em}{}

\titleformat{\subsection}
{\large\bfseries\color{bleufonce}}
{}{0em}{}

% --------------------------------------------------
% Environnements
% --------------------------------------------------
\newtcolorbox{consignesbox}{
    colback=bleuclair,
    colframe=bleufonce,
    boxrule=0.8pt,
    arc=2mm,
    left=2mm,right=2mm,top=1.5mm,bottom=1.5mm
}

\newtcolorbox{espritbox}{
    colback=grisclair,
    colframe=black!60,
    boxrule=0.6pt,
    arc=2mm,
    left=2mm,right=2mm,top=1.5mm,bottom=1.5mm
}

\newtcolorbox{exobox}[1]{
    colback=white,
    colframe=bleufonce,
    boxrule=1pt,
    arc=1.8mm,
    title=\textbf{#1},
    coltitle=white,
    colbacktitle=bleufonce,
    fonttitle=\bfseries,
    left=2mm,right=2mm,top=1.5mm,bottom=1.5mm
}

\newtcolorbox{objectifbox}{
    colback=bleuclair,
    colframe=bleufonce,
    boxrule=0.5pt,
    arc=1mm,
    left=1.5mm,right=1.5mm,top=1mm,bottom=1mm
}

% --------------------------------------------------
% Commandes utiles
% --------------------------------------------------
\newcommand{\vect}[1]{\overrightarrow{#1}}
\newcommand{\R}{\mathbb{R}}
\newcommand{\e}{\mathrm{e}}
\newcommand{\dd}{\,\mathrm{d}}

% --------------------------------------------------
% Document
% --------------------------------------------------
\begin{document}

\begin{center}
    {\Huge \bfseries \color{bleufonce} Devoir Maison de Physique-Chimie}\\[0.4em]
    {\Large Niveau Première générale}\\[0.4em]
    \rule{0.7\linewidth}{0.6pt}\\[0.4em]
    Durée conseillée : \textbf{4 h}
\end{center}

\begin{consignesbox}
\textbf{Consignes pour le DM :}
\begin{itemize}[label=---]
    \item proposer une rédaction de type DS, soignée, structurée et claire ;
    \item faire tous les exercices ;
    \item toute affirmation doit être justifiée ;
    \item lorsqu’une question est qualitative, on attend un raisonnement précis ;
    \item noter dans la marge le temps consacré à chaque exercice ;
    \item lorsqu’un schéma ou un tableau aide le raisonnement, il est conseillé d’en faire un.
\end{itemize}
\end{consignesbox}

\begin{espritbox}
\textbf{Esprit du sujet.} Ce devoir ne cherche pas seulement à vérifier des automatismes. Il vise à tester la compréhension des notions de Première, la capacité à relier plusieurs idées du cours, à discuter un modèle, à interpréter une expérience et à rédiger une réponse argumentée.
\end{espritbox}

% ==================================================
\begin{exobox}{Exercice 1 -- Dissolution, dilution et précision expérimentale}
\begin{objectifbox}
\textbf{Objectif :} comprendre précisément ce que l’on fait lorsqu’on prépare une solution, distinguer dissolution et dilution, et réfléchir à la qualité d’un protocole.
\end{objectifbox}

On souhaite préparer une solution aqueuse de sulfate de cuivre, puis en réaliser une dilution.

On dispose du matériel suivant :
\begin{itemize}
    \item balance électronique ;
    \item coupelle de pesée et spatule ;
    \item bécher ;
    \item entonnoir ;
    \item fiole jaugée de \(100,0~\mathrm{mL}\) ;
    \item pipette jaugée de \(10,0~\mathrm{mL}\) ;
    \item propipette ;
    \item pissette d’eau distillée.
\end{itemize}

\begin{center}
\begin{tikzpicture}[scale=1]
% balance
\draw[fill=gray!15] (-5,0) rectangle (-2.8,1.2);
\draw[fill=gray!5] (-4.55,0.65) rectangle (-3.25,1.05);
\node at (-3.9,0.3) {\small balance};

% bécher
\draw[thick] (-0.9,0) -- (-0.9,1.5);
\draw[thick] (0.9,0) -- (0.9,1.5);
\draw[thick] (-0.9,0) -- (0.9,0);
\draw[thick] (-0.9,1.5) arc[start angle=180,end angle=0,x radius=0.9,y radius=0.15];
\fill[blue!20] (-0.88,0.02) rectangle (0.88,0.75);
\node at (0,-0.35) {\small bécher};

% fiole
\draw[thick] (3.4,0) .. controls (3.1,0.9) and (2.8,1.6) .. (3.5,2.2);
\draw[thick] (4.6,0) .. controls (4.9,0.9) and (5.2,1.6) .. (4.5,2.2);
\draw[thick] (3.5,2.2) -- (3.5,3.8);
\draw[thick] (4.5,2.2) -- (4.5,3.8);
\draw[thick] (3.5,3.8) -- (4.5,3.8);
\draw[thick,red] (3.45,3.1) -- (4.55,3.1);
\fill[blue!15] (3.52,0.12) .. controls (3.35,0.85) and (3.2,1.4) .. (3.6,2.0)
-- (4.4,2.0) .. controls (4.8,1.4) and (4.65,0.85) .. (4.48,0.12) -- cycle;
\node at (4,-0.35) {\small fiole jaugée};
\node[red] at (5.15,3.1) {\small trait de jauge};

\end{tikzpicture}
\end{center}

\begin{enumerate}[label=\textbf{\arabic*.}]
    \item Expliquer clairement la différence entre une \textbf{dissolution} et une \textbf{dilution}. On attend plus qu’une simple phrase : il faut préciser ce qui change et ce qui ne change pas dans chaque opération.
    
    \item Définir la \textbf{concentration en masse} \(C_m\) d’une solution et donner sa formule.

    \item On dissout une masse \(m\) de soluté dans un volume \(V\) de solution. Établir la relation donnant \(C_m\) en fonction de \(m\) et \(V\).

    \item Décrire précisément un protocole expérimental permettant de préparer \(100,0~\mathrm{mL}\) d’une solution par dissolution d’un solide. Le vocabulaire expérimental doit être rigoureux.

    \item Pourquoi est-il important de dissoudre d’abord le solide dans une petite quantité d’eau avant d’ajuster au trait de jauge ?

    \item On prélève ensuite \(10,0~\mathrm{mL}\) de la solution mère pour préparer \(100,0~\mathrm{mL}\) de solution fille.
    \begin{enumerate}[label=\alph*)]
        \item Déterminer le facteur de dilution.
        \item Expliquer précisément ce que signifie physiquement ce facteur.
    \end{enumerate}

    \item Un élève dit : \og dans une dilution, on ajoute de l’eau, donc on ajoute de la matière au mélange, donc la quantité de soluté augmente \fg.  
    Discuter cette affirmation avec soin.

    \item On hésite entre utiliser une pipette jaugée et un simple bécher gradué pour prélever la solution mère. Quel choix est le plus pertinent ? Justifier.

    \item Question de compréhension.  
    Deux élèves préparent deux solutions filles différentes à partir de la même solution mère :
    \begin{itemize}
        \item le premier prélève \(10,0~\mathrm{mL}\) puis complète à \(100,0~\mathrm{mL}\) ;
        \item le second prélève \(20,0~\mathrm{mL}\) puis complète à \(100,0~\mathrm{mL}\).
    \end{itemize}
    Sans calcul numérique compliqué, comparer les concentrations finales et expliquer pourquoi.

    \item Question plus fine.  
    On veut préparer une solution fille très peu concentrée.  
    Vaut-il mieux faire directement une très forte dilution en une seule étape, ou bien plusieurs dilutions successives ?  
    On attend une discussion argumentée portant à la fois sur la \textbf{théorie} et sur la \textbf{précision expérimentale}.

    \item \textbf{Problème ouvert.}  
    Proposer une méthode expérimentale réaliste pour vérifier qu’une dilution a bien été réalisée correctement, sans connaître à l’avance la concentration exacte de la solution.
\end{enumerate}
\end{exobox}

% ==================================================
\begin{exobox}{Exercice 2 -- Couleur, absorption et interprétation physique}
\begin{objectifbox}
\textbf{Objectif :} dépasser la simple récitation du cours sur les couleurs et relier couleur perçue, lumière transmise, lumière absorbée et conditions d’observation.
\end{objectifbox}

On étudie trois solutions transparentes :
\[
\text{A : bleue} \qquad \text{B : rouge} \qquad \text{C : verte}
\]

On rappelle qu’un objet ou une solution colorée n’émet pas forcément lui-même de lumière visible : il modifie la lumière qu’il reçoit.

\begin{center}
\begin{tikzpicture}[scale=1]
% source
\draw[fill=yellow!30,thick] (-5,0) circle (0.45);
\node at (-5,0) {\small source};

% faisceau incident
\draw[very thick,->] (-4.5,0) -- (-2.1,0);

% cuve
\draw[thick] (-2,-1.2) rectangle (0,1.2);
\fill[blue!20] (-2,-1.2) rectangle (0,1.2);
\node at (-1,-1.55) {\small solution};

% sortie
\draw[very thick,->,blue] (0.1,0.2) -- (2.7,0.2);
\draw[very thick,->,red!70!black] (0.1,-0.2) -- (2.2,-0.2);
\draw[very thick,->,green!60!black] (0.1,0) -- (2.45,0);

\node at (3.5,0.8) {\small lumière transmise};
\node at (3.3,-0.9) {\small certaines radiations};
\node at (3.25,-1.3) {\small peuvent être absorbées};

\end{tikzpicture}
\end{center}

\begin{enumerate}[label=\textbf{\arabic*.}]
    \item Qu’appelle-t-on \textbf{lumière blanche} ?

    \item Expliquer, avec une phrase scientifiquement correcte, pourquoi une solution peut paraître colorée lorsqu’elle est éclairée en lumière blanche.

    \item Une solution paraît bleue. Quelle information cela donne-t-il sur la lumière qu’elle transmet principalement ?

    \item Cette même solution absorbe-t-elle plutôt les radiations bleues, rouges, vertes, ou un mélange de plusieurs radiations ? Discuter qualitativement.

    \item On éclaire maintenant une solution rouge avec une lumière uniquement verte.
    \begin{enumerate}[label=\alph*)]
        \item Que peut-on prévoir pour son aspect ?
        \item Pourquoi ?
    \end{enumerate}

    \item Même question pour une solution verte éclairée uniquement en lumière rouge.

    \item Un élève affirme : \og une solution bleue contient de la lumière bleue \fg.  
    Cette phrase est-elle physiquement satisfaisante ? La reformuler correctement.

    \item Une solution noire absorbe presque toute la lumière visible incidente.  
    Une solution transparente incolore n’absorbe presque aucune radiation visible.  
    Comparer soigneusement ces deux situations.

    \item Question de compréhension plus poussée.  
    Une boisson paraît rouge sous lumière blanche, mais presque noire sous un éclairage bleu.  
    Expliquer ce phénomène sans contradiction.

    \item On place successivement un filtre rouge, puis un filtre vert, puis un filtre bleu devant une solution verte. Décrire qualitativement ce qu’on peut attendre dans chaque cas.

    \item Pourquoi l’observation de la couleur seule ne permet-elle pas toujours d’identifier avec certitude une espèce chimique ?

    \item \textbf{Problème ouvert.}  
    On veut comparer expérimentalement deux solutions colorées pour savoir laquelle absorbe le plus la lumière rouge.  
    Proposer un protocole simple de laboratoire, expliquer le principe physique utilisé, et préciser quelles observations permettraient de conclure.
\end{enumerate}
\end{exobox}

% ==================================================
\begin{exobox}{Exercice 3 -- Mouvement, forces et discussion d’un modèle}
\begin{objectifbox}
\textbf{Objectif :} comprendre profondément le principe d’inertie, le lien entre forces et mouvement, et éviter les erreurs classiques de raisonnement.
\end{objectifbox}

On étudie le mouvement d’un palet sur une table horizontale.

\begin{center}
\begin{tikzpicture}[scale=1.05]
% table
\draw[thick] (-5,0) -- (5.5,0);

% palet
\draw[fill=gray!25,thick] (-1,0.12) rectangle (-0.1,0.62);
\node at (-0.55,0.37) {\small palet};

% vitesse
\draw[very thick,->] (-0.05,0.37) -- (2.2,0.37);
\node at (1.1,0.7) {\(\vec v\)};

% poids
\draw[very thick,->] (-0.55,0.12) -- (-0.55,-1.3);
\node[left] at (-0.55,-0.7) {\(\vec P\)};

% réaction
\draw[very thick,->] (-0.55,0.62) -- (-0.55,1.8);
\node[left] at (-0.55,1.25) {\(\vec R\)};

% frottement éventuel
\draw[very thick,->,red!70!black] (-0.1,0.37) -- (-1.4,0.37);
\node[red!70!black] at (-1.1,0.8) {\(\vec f\) éventuel};

\end{tikzpicture}
\end{center}

Dans un premier temps, les frottements sont supposés négligeables.

\begin{enumerate}[label=\textbf{\arabic*.}]
    \item Dans le référentiel du laboratoire, qu’appelle-t-on un \textbf{mouvement rectiligne uniforme} ?

    \item Quelles sont les forces qui s’exercent sur le palet après qu’il a été lancé, si les frottements sont négligeables ?

    \item Représenter ces forces et préciser leurs directions.

    \item Expliquer pourquoi le principe d’inertie permet alors de prévoir un mouvement rectiligne uniforme.

    \item Un élève affirme : \og s’il n’y a plus de force dans le sens du mouvement, le mobile doit s’arrêter \fg.  
    Dire si cette affirmation est vraie ou fausse, puis la corriger rigoureusement.

    \item Pourquoi confond-on souvent, à tort, les idées de \og mouvement \fg\ et de \og force \fg\ ?

    \item Dans une seconde expérience, le palet ralentit progressivement.
    \begin{enumerate}[label=\alph*)]
        \item Qu’est-ce que cela indique sur la résultante des forces ?
        \item Quelle force peut expliquer cette observation ?
    \end{enumerate}

    \item Peut-on avoir :
    \begin{enumerate}[label=\alph*)]
        \item une vitesse non nulle et une résultante des forces nulle ?
        \item une vitesse nulle et une résultante des forces non nulle ?
    \end{enumerate}
    Dans chaque cas, justifier par un exemple ou un contre-exemple.

    \item Deux mobiles différents ont à un instant donné la même vitesse.  
    Peut-on en conclure qu’ils subissent les mêmes forces ? Expliquer.

    \item Deux mobiles différents subissent la même résultante des forces.  
    Peut-on en conclure qu’ils ont le même mouvement ? Expliquer qualitativement.

    \item Question plus subtile.  
    Un palet se déplace en ligne droite, mais sa vitesse augmente.  
    Que peut-on dire de la direction de la résultante des forces ?

    \item \textbf{Problème ouvert.}  
    Dans quels cas le modèle \og frottements négligeables \fg\ est-il raisonnable ?  
    Dans quels cas devient-il insuffisant ?  
    On attend une discussion argumentée, avec exemples concrets, portant sur les effets possibles des frottements sur le mouvement.
\end{enumerate}
\end{exobox}

% ==================================================
\begin{exobox}{Exercice 4 -- Énergie, transferts et puissance}
\begin{objectifbox}
\textbf{Objectif :} distinguer clairement énergie et puissance, et raisonner sur des transferts énergétiques dans des situations concrètes.
\end{objectifbox}

Une élève de masse \(m\) monte un escalier pour atteindre un étage situé à une hauteur \(h\).

\begin{center}
\begin{tikzpicture}[scale=0.95]
% escalier
\draw[thick] (-4,0) -- (-2.8,0) -- (-2.8,0.6) -- (-1.6,0.6) -- (-1.6,1.2) -- (-0.4,1.2)
-- (-0.4,1.8) -- (0.8,1.8) -- (0.8,2.4) -- (2,2.4);

% personnage bas
\draw[thick] (-3.4,0.7) circle (0.16);
\draw[thick] (-3.4,0.54) -- (-3.4,0.05);
\draw[thick] (-3.4,0.38) -- (-3.7,0.18);
\draw[thick] (-3.4,0.38) -- (-3.1,0.18);
\draw[thick] (-3.4,0.05) -- (-3.6,-0.3);
\draw[thick] (-3.4,0.05) -- (-3.15,-0.28);

% personnage haut
\draw[thick] (1.5,3.1) circle (0.16);
\draw[thick] (1.5,2.94) -- (1.5,2.45);
\draw[thick] (1.5,2.78) -- (1.2,2.58);
\draw[thick] (1.5,2.78) -- (1.8,2.58);
\draw[thick] (1.5,2.45) -- (1.3,2.1);
\draw[thick] (1.5,2.45) -- (1.75,2.15);

% fleche hauteur
\draw[very thick,<->,blue] (2.7,0) -- (2.7,3.0);
\node[blue,right] at (2.7,1.5) {\(h\)};

\end{tikzpicture}
\end{center}

\begin{enumerate}[label=\textbf{\arabic*.}]
    \item Rappeler l’expression de l’énergie potentielle de pesanteur.

    \item Lorsque l’élève monte, cette énergie augmente-t-elle ou diminue-t-elle ? Pourquoi ?

    \item Exprimer la variation d’énergie potentielle de pesanteur entre le bas et le haut de l’escalier.

    \item Interpréter physiquement cette variation.

    \item Deux élèves de même masse montent au même étage :
    \begin{itemize}
        \item l’un monte lentement ;
        \item l’autre monte rapidement.
    \end{itemize}
    Comparer :
    \begin{enumerate}[label=\alph*)]
        \item la variation d’énergie potentielle ;
        \item la puissance moyenne développée.
    \end{enumerate}

    \item Définir la puissance moyenne.

    \item Expliquer pourquoi deux actions peuvent correspondre au même transfert d’énergie mais à des puissances différentes.

    \item Un élève affirme : \og celui qui fournit la plus grande puissance fournit forcément la plus grande énergie \fg.  
    Discuter soigneusement cette affirmation.

    \item On compare maintenant :
    \begin{itemize}
        \item la montée de l’escalier à pied ;
        \item la montée par ascenseur.
    \end{itemize}
    Dans les deux cas, l’altitude finale est la même.
    \begin{enumerate}[label=\alph*)]
        \item Quelle grandeur liée à la pesanteur est forcément la même à l’arrivée ?
        \item Quelles différences peut-on néanmoins observer du point de vue énergétique ?
    \end{enumerate}

    \item Pourquoi l’énergie fournie par un système n’est-elle pas toujours entièrement convertie sous la forme utile attendue ?

    \item Donner deux exemples concrets de \textbf{pertes} ou de \textbf{dissipations} dans des situations de la vie courante.

    \item \textbf{Problème ouvert.}  
    Rédiger une discussion argumentée sur les transferts d’énergie mis en jeu lorsqu’une personne monte un escalier.  
    On attend notamment :
    \begin{itemize}
        \item une identification des systèmes concernés ;
        \item les principales formes d’énergie en jeu ;
        \item une réflexion sur l’utilité de distinguer énergie et puissance ;
        \item une conclusion claire.
    \end{itemize}
\end{enumerate}
\end{exobox}

% ==================================================
\begin{exobox}{Exercice 5 -- Ondes sonores, période, fréquence et interprétation}
\begin{objectifbox}
\textbf{Objectif :} relier la représentation temporelle d’un signal à la notion physique de son, et éviter les confusions classiques sur la propagation d’une onde.
\end{objectifbox}

On étudie le signal sonore émis par un diapason. Sur un enregistrement temporel, on observe la répétition régulière d’un même motif.

\begin{center}
\begin{tikzpicture}[scale=1]
\draw[->] (0,0) -- (12.5,0) node[right] {temps};
\draw[->] (0,-2) -- (0,2.2) node[above] {signal};

\draw[domain=0:12,samples=300,smooth,thick,blue]
plot (\x,{1.2*sin(deg(2*pi*\x/2.4))});

\draw[dashed] (0,0) -- (12,0);

\draw[decorate,decoration={brace,mirror,amplitude=5pt}] (1.2,-1.55) -- (3.6,-1.55);
\node at (2.4,-1.95) {\(T\)};

\draw[dashed] (1.2,-1.4) -- (1.2,1.2);
\draw[dashed] (3.6,-1.4) -- (3.6,1.2);
\end{tikzpicture}
\end{center}

On rappelle que
\[
f=\frac{1}{T}.
\]

\begin{enumerate}[label=\textbf{\arabic*.}]
    \item Définir ce qu’est un \textbf{signal périodique}.

    \item Définir la \textbf{période} \(T\) d’un signal.

    \item Définir la \textbf{fréquence} \(f\) d’un signal et donner son unité.

    \item Un enregistrement montre que \(5\) périodes occupent une durée totale de \(0{,}020~\mathrm{s}\).
    \begin{enumerate}[label=\alph*)]
        \item Déterminer la période \(T\).
        \item En déduire la fréquence \(f\).
    \end{enumerate}

    \item Si la fréquence augmente, le son est-il perçu comme plus grave ou plus aigu ?

    \item Deux sons ont la même fréquence mais des signaux de formes différentes.  
    Que peut-on dire :
    \begin{enumerate}[label=\alph*)]
        \item de leur hauteur ;
        \item de leur timbre ?
    \end{enumerate}

    \item Un élève affirme : \og le son se déplace parce que l’air se déplace globalement de la source vers l’auditeur \fg.  
    Expliquer pourquoi cette formulation est trompeuse.

    \item Pourquoi le son ne se propage-t-il pas dans le vide ?

    \item Un son fort se propage-t-il plus vite qu’un son faible ? Justifier.

    \item Peut-on entendre un son derrière un obstacle ? Donner une réponse nuancée.

    \item Question plus fine.  
    Pourquoi la représentation temporelle d’un signal enregistrée par un microphone ne correspond-elle pas à \og la forme visible du son dans l’air \fg, même si elle en donne une information utile ?

    \item \textbf{Problème ouvert.}  
    Proposer une méthode expérimentale permettant de mesurer la fréquence d’un son à l’aide d’un microphone et d’un enregistrement temporel.  
    On attend :
    \begin{itemize}
        \item les mesures à effectuer ;
        \item la manière d’améliorer la précision ;
        \item la formule finale utilisée ;
        \item une remarque sur une source possible d’erreur.
    \end{itemize}
\end{enumerate}
\end{exobox}

\vfill
\begin{center}
    \rule{0.5\linewidth}{0.5pt}\\
    \textit{Fin du devoir}
\end{center}

\end{document}